% Original topic-level problem, solution and MCQ reasoning.
\subsection{বাস্তব সমস্যা: চলন্ত Bus-এ Call ও Data}
\chTwoPair{\textbf{সমস্যা:} বাস চলার সময় ফোনে কথা চলছে। পরে অন্য অপারেটরের সহযোগী নেটওয়ার্কে সেবা নিতে হলো। একই ফোনে ভিডিও কখনো দ্রুত চলে, কখনো আটকে যায়।\par
\textbf{সমাধান:} সেবাদানকারী সেল বদলানো হলো Handoff। নিজস্ব নেটওয়ার্কের বাইরে সহযোগীর সেবা নেওয়া হলো Roaming। ভিডিও চলার মান সংকেত, বহনক্ষমতা, ভিড় ও বিলম্ব অনুযায়ী বদলায়। শুধু 4G লেখা থাকলেই নির্দিষ্ট ডেটা হার নিশ্চিত হয় না।}{\textbf{Mobile Link-এর Uplink ও Downlink:} ফোন থেকে স্টেশনে পাঠানো হলো Uplink; স্টেশন থেকে ফোনে পাঠানো হলো Downlink। স্যাটেলাইটে ভূকেন্দ্র থেকে উপগ্রহের নামের সঙ্গে কোন প্রান্তকে ভিত্তি ধরা হচ্ছে, তা মিলিয়ে বুঝতে হয়।\par
\textbf{Packet Data ও Generation:} GPRS/EDGE প্রচলিত 2G-এর ডেটা সেবার উন্নত রূপ। প্যাকেট ডেটা একটিমাত্র প্রজন্মে সীমাবদ্ধ নয়।\par}
\begin{multicols}{2}
\textbf{SMS, MMS ও Billing:} SMS-এর পূর্ণরূপ Short Message Service এবং MMS-এর পূর্ণরূপ Multimedia Messaging Service। SMS দিয়ে লেখা এবং MMS দিয়ে ছবি বা অডিওসহ বার্তা পাঠানো যায়। প্রচলিত প্রশ্নে এগুলো ডিজিটাল বা 2G ধারার পরিচয় হিসেবে থাকতে পারে। Prepaid হলো আগে অর্থ দিয়ে সেবা ব্যবহারের বিল পদ্ধতি। প্রচলিত প্রজন্মের ছকে এটি 2G-এর সঙ্গে দেখানো হয়।\par
\textbf{Circuit ও Packet Switching:} Circuit Switching-এ একটি যোগাযোগ পর্বের জন্য নির্দিষ্ট যোগাযোগসম্পদ সংরক্ষণ করা হয়। Packet Switching-এ ডেটা ছোট প্যাকেটে ভাগ করে নেটওয়ার্কের সম্পদ ভাগাভাগি করা হয়। GPRS-এও প্যাকেট ডেটা আছে। 3G-তে এর ব্যবহার বিস্তৃত হয় এবং 4G-তে IP-ভিত্তিক ব্যবস্থা প্রধান।\par
\textbf{Green Phone-এর পরিচয়:} পরীক্ষায় Green Phone নামটি CDMA-এর প্রচলিত পরিচয় হিসেবে থাকতে পারে। CDMA-র শক্তি নিয়ন্ত্রণ প্রয়োজনমতো প্রেরণশক্তি রেখে বিঘ্ন ও ব্যাটারির ব্যবহার কমাতে সাহায্য করে। Green Phone আধুনিক মান বা প্রত্যয়নপত্রের নাম নয়।\par
\textbf{Portable Cellular Phone-এর ইতিহাস:} Motorola-র Martin Cooper ৩ এপ্রিল ১৯৭৩ সালে হাতে বহনযোগ্য সেলুলার ফোনে প্রকাশ্যে কল করেন। DynaTAC 8000X বাণিজ্যিক বহনযোগ্য ফোন হিসেবে আসে ১৯৮৩ সালে। প্রথম প্রদর্শন ও বাজারে আসার সাল পৃথক। Cooper-এর অবদান সেলুলার হ্যান্ডসেটের ইতিহাসে গুরুত্বপূর্ণ; তিনি সব বেতার যোগাযোগের একমাত্র উদ্ভাবক নন।\par
\end{multicols}
\textbf{পরিষেবার Rate:} VoIP, Streaming Video ও HD Video Conference-এর প্রচলিত সংখ্যাগত উদাহরণ যোগাযোগ ও ব্যান্ডউইথের অংশে দেওয়া আছে (অংশ \ref{ch2:c01-12}); সেখানকার পরিষেবা-ছক অনুসরণ করবে।

\textbf{অবস্থান ও ডেটা সেবা:} GPS = Global Positioning System। এটি অবস্থান নির্ণয়ের স্যাটেলাইট ব্যবস্থা। GPRS-এর মোবাইল ডেটা সেবা থেকে GPS আলাদা। IP টেলিফোনি, গেমিং এবং HD/3D ভিডিও মোবাইল ব্রডব্যান্ডের ব্যবহার। এসব কাজে প্রয়োজনীয় ডেটা হার, যন্ত্র ও কোডেক সামঞ্জস্যপূর্ণ হতে হয়।
